\section{Related Work}
\label{sec:related_work}

\paragraph{Dynamic manipulation.}
Compensating for the time between sensing and acting is a classical problem, from
control with dead time \cite{smith1957} to catching systems that track the target
with Kalman-type filters and plan an interception
\cite{barshalom2001,allen1993,baeuml2010,marturi2019dynamic}; such systems are
precise but engineered for each object and task. Learned policies have recently been
brought to moving scenes \cite{dynamicvla}. PUMA reads historical optical flow
and is supervised with future object features \cite{domino}, and DynamicWAM conditions
a WAM on flow and kinematic descriptors and executes it with real-time chunking
\cite{dynamicwam}. Both infer past motion from image-space cues, whereas \ours takes
the target's metric SE(3) pose history as input and observes only the current RGB-D
frame. Other dynamic policies forecast salient tokens \cite{ahead}, future images
\cite{dbctfp}, compact future codes \cite{foretime}, or 3D scene dynamics
\cite{physmani} and condition their actions on these forecasts at test time. DynaWM
adapts a base VLA's action chunks for dynamic manipulation \cite{dynawm}, and
DynamicManip synthesizes dynamic demonstrations from a static one \cite{dynamicmanip}.
Execution schemes reduce or hide latency by overlapping inference with execution
\cite{dynamicvla,smolvla}, keeping consecutive chunks consistent \cite{rtc,bid}, or
conditioning on the execution time \cite{vlash}. They decide when actions are computed
and applied, not what the policy represents. DOMINO \cite{domino} and PhysMani-Bench
\cite{physmani} provide diverse dynamic-manipulation tasks, but their released default
evaluation pipelines advance the simulation only after a policy has returned its
prediction. \bench advances the scene during inference and charges each method its
measured latency.

\paragraph{Predictive supervision in WAMs.}
World--action models (WAMs) generate future video and actions in one network
\cite{uwm,uva,dreamzero,openwam_alpha}, building on video prediction as a policy
objective \cite{gr1,seer}. Because generating video is costly, recent work changes what
is forecast, such as dynamic regions, depth, and semantic features \cite{dreamvla}, flow
and geometry alongside RGB \cite{dreamwam}, object masks or slots \cite{maskwam,oawam},
or compact latents \cite{flare,wog}, and asks whether the forecast is needed at test
time. Fast-WAM retains most of the benefit of video modeling without test-time
imagination \cite{fastwam}, and PUMA and SG-WAM drop their predictors after training
\cite{domino,sgwam}. \ours keeps this training-only use but forecasts primarily a metric
4D state rather than pixels, features, or latents, so the target's future pose is
supervised explicitly instead of being left implicit.

\paragraph{Object-centric and geometric motion representations.}
Metric 3D inputs make policies more precise and sample-efficient \cite{peract,dp3}.
Motion has been represented by 2D point tracks \cite{atm,track2act,im2flow2act}, 3D
flow \cite{generalflow,flowdreamer,pointworld,physmani}, and object keypoints or SE(3)
trajectories \cite{kpam,spot,objectforesight}, while rigid-body decompositions have long
been used to learn scene dynamics \cite{se3nets,se3posenets}. These works describe the
scene with one kind of representation, such as dense flow or tracks everywhere, or a
pose for each object. Target-centric predictors, including dynamic VLAs with
object-level forecasts, model only the target's future \cite{domino,spot} and leave out
the bodies that block or push it. DOCIR separates the robot, the target, and obstacles,
but uses these distinctions only as input encodings \cite{docir}. \ours assigns
precision by role: a coarse sparse 3D flow describes the whole scene, while the target,
on which contact depends to within millimeters, is additionally described by its SE(3)
history and forecast.
